% English version of paper/ja_v2/08_限界.md (2026-11-17), condensed; restored items after the paper-prose review (reviews/v2_en/prose.md 2-1 to 2-9).
\section{Limitations}\label{sec:limitations}

\textbf{Scope.} RQ1(a) uses the single-source and all-source settings, and RQ1(b) and RQ2 the LOCO setting (four corpora; RQ2 adds CREMA-D), so Sect.~\ref{sec:discussion} connects results from different settings. All corpora are in English, IEMOCAP in the LOCO setting is only its improvised part, and one encoder per modality was used (plus HuBERT in the single-source setting). The results concern linear classifiers on frozen features. Fine-tuned encoders have been reported to give better within-corpus dimensional emotion prediction (CCC, the concordance correlation coefficient) than frozen features \cite{wagner2023}, and DANN and MME added their loss terms to a linear map and classifier instead of training a feature extractor. The results should therefore be read as properties of this diagnostic setting; testing other corpora, encoders and fine-tuning is the next step.

\textbf{Answer bias and the Gaussian approximation.} The findings that mean subtraction almost removed answer bias and that the K labels mainly corrected bias come from the single-source setting, with means over all target utterances including test utterances. In the LOCO setting, an intercept-only fix on the K labels was below the K-label baseline in point estimates in 18 of 32 conditions, so which error the K labels correct there is unknown. Direction losses rely on an equal-variance Gaussian approximation that failed the pre-set criterion for visual features (median approximation error 0.47 times the median direction loss); the visual $d_T/D^{*}$ in Sect.~\ref{sec:rq1a} and the visual rows of Table~\ref{tab:allsource} are values of these features and should be read with this caveat.

\textbf{Analyses and methods built on the same data.} The advantage analysis is descriptive, set up after part of the saved outputs had been seen, and its counts and correlations are not tests. The forms, grids and selection procedures of the methods were built on the evaluation corpora, whose test folds also served as provisional-target tests for other targets. The idea of the proposal followed the evaluation run of Sect.~\ref{sec:rq1b}, CREMA-D was added after its provisional values had been seen, and UARs and differences were inspected during and after that evaluation run. Differences from the product and other combinations may therefore be biased in favour of the proposal. The strongest-single-classifier comparison was added afterwards from saved outputs, covers estimation-time decisions only and was not computed for CREMA-D. Evaluating once on a new corpus with the procedure fixed is the next step.

\textbf{Retrained classifiers and pseudo-labels.} Only AVT (AV for CREMA-D) was retrained on the proposal's pseudo-labels, not the audio classifier that gained most from distillation; retraining the audio, visual and text classifiers on the proposal's pseudo-labels and comparing them with the product is the next step. The pseudo-labels include the retrained classifier's own modality. Section~\ref{sec:rq1b-margin} also lists providers without the own modality, but the contribution of the other modalities to the product was not separated, so the gains cannot be attributed to other modalities alone.

\textbf{Shared content and reversed speech.} Speakers are separated, but the same person (MOSEI), episode (MELD) or storyline (IEMOCAP) may appear on both sides, and in CREMA-D the same sentences appear on both sides because actors read 12 fixed sentences. Reversed speech was checked with Whisper transcripts, not with human listeners. Reversal also inverts time-directional cues such as intonation contours and produces speech unlike the data WavLM was trained on, so part of the drop in $D$ may come from this change rather than from the loss of words.

\textbf{Advantage on a new corpus.} The advantage was measured with all candidate labels and only at K = 5 and 25, and a new corpus lacks these labels. Estimating the advantage within the K labels and comparing it with the value on all candidate labels is the next step.

\textbf{Intervals.} Intervals are not corrected for multiple comparisons, and counts of corpora or conditions are not tests.
